\documentclass[11pt]{article}
\usepackage[spanish,es-nodecimaldot]{babel}
\usepackage{amsmath,amssymb}
\usepackage{geometry}
\usepackage{microtype}
\usepackage{array}
\usepackage{booktabs}
\usepackage{xcolor}
\usepackage{hyperref}

\geometry{letterpaper,margin=2.5cm}

\hypersetup{
  colorlinks=true,
  linkcolor=blue!55!black,
  urlcolor=blue!55!black,
  citecolor=blue!55!black
}

\setlength{\parindent}{0pt}
\setlength{\parskip}{0.65em}
\setlength{\jot}{0.4em}

\newcommand{\curso}{F\'isica Aplicada a las Ciencias Farmac\'euticas}
\newcommand{\unidad}{Elementos de mec\'anica de la part\'icula en una dimensi\'on}
\newcommand{\vect}[1]{\boldsymbol{#1}}
\newcommand{\uvect}[1]{\hat{\boldsymbol{#1}}}
\newcommand{\R}{\mathbb{R}}
\newcommand{\dd}{\,\mathrm{d}}
\newcommand{\unidadSI}[1]{\,\mathrm{#1}}
\newcommand{\clase}[1]{\section*{#1}\addcontentsline{toc}{section}{#1}}
\newcommand{\subclase}[1]{\subsection*{#1}\addcontentsline{toc}{subsection}{#1}}

\newenvironment{objetivos}
  {\subclase{Objetivos de la clase}\begin{itemize}}
  {\end{itemize}}

\newenvironment{pizarra}
  {\subclase{Para desarrollar en pizarra}\begin{enumerate}}
  {\end{enumerate}}

\newenvironment{ejemplo}[1]
  {\subclase{Ejemplo: #1}}
  {}


\begin{document}

\portada{1}{Medir, modelar y estimar}{Unidad I: Elementos de mecánica de la partícula en una dimensión}

\pagenumbering{roman}
\tableofcontents
\inicioContenido

\section{Mapa conceptual}

\mapalineal{
  \node[concepto] (fen) {fenómeno};
  \node[concepto,right=of fen] (mod) {modelo};
  \node[concepto,right=of mod] (mag) {magnitudes};
  \node[concepto,right=of mag] (uni) {unidades};
  \node[concepto,right=of uni] (ecu) {relaciones};
  \node[concepto,right=of ecu] (con) {control físico};
  \draw[flecha] (fen)--(mod);
  \draw[flecha] (mod)--(mag);
  \draw[flecha] (mag)--(uni);
  \draw[flecha] (uni)--(ecu);
  \draw[flecha] (ecu)--(con);
}

Una descripción física comienza con un fenómeno, pero no intenta copiar todos
sus detalles. Se escoge un modelo, se identifican las magnitudes relevantes y
se expresan con unidades. Las ecuaciones relacionan esas magnitudes. Antes de
aceptar un resultado se revisan sus unidades, sus dimensiones y su escala.

\section{Medir exige una pregunta}

Una botella con líquido admite muchas preguntas: ¿cuánta masa contiene?,
¿qué volumen ocupa?, ¿a qué temperatura está?, ¿cuánto tarda en vaciarse? Cada
pregunta selecciona propiedades distintas. La masa puede ser importante para
preparar una formulación, mientras que el tiempo de vaciado puede ser
irrelevante para esa misma tarea.

\begin{center}
\begin{tikzpicture}[
  objeto/.style={draw=GrisTexto,thick,rounded corners=1mm},
  rotulo/.style={font=\small,align=center},
  flecha/.style={-{Stealth[length=2mm]},draw=GrisLinea}
]
  \draw[objeto] (0,0) rectangle (2.1,3.0);
  \draw[AzulTitulo,thick] (0.15,1.7) -- (1.95,1.7);
  \node[rotulo] at (1.05,0.9) {líquido};
  \node[rotulo] (real) at (1.05,-0.45) {sistema real};

  \node[rotulo,draw=GrisLinea,rounded corners=1mm,minimum width=3.1cm,
        minimum height=1.2cm] (modelo) at (5.0,1.5)
        {modelo homogéneo\\\(\rho=\text{constante}\)};
  \draw[flecha] (2.3,1.5)--(3.35,1.5);
  \node[font=\small] at (2.83,1.82) {simplificar};

  \node[rotulo,draw=GrisLinea,rounded corners=1mm,minimum width=3.0cm,
        minimum height=1.2cm] (datos) at (8.9,1.5)
        {\(m,\ V,\ \rho\)\\magnitudes};
  \draw[flecha] (6.65,1.5)--(7.35,1.5);
\end{tikzpicture}
\end{center}

Un \textbf{sistema} es la parte de la realidad que se estudia. Un
\textbf{modelo} es una representación deliberadamente simplificada del
sistema. Si el líquido se considera homogéneo, su densidad puede tratarse como
constante; si hay sedimento o gradientes de temperatura, ese supuesto puede
fallar. Un modelo no es verdadero en cualquier circunstancia: es útil dentro
de un dominio.

\textbf{Medir} significa comparar una propiedad con un patrón acordado. Una
medición se expresa como
\[
  \text{cantidad física}
  =
  \text{valor numérico}\times\text{unidad}.
\]
El número sin unidad está incompleto. \(5\) puede representar cinco metros,
cinco segundos o cinco miligramos; cada opción describe una realidad distinta.

\begin{ejemplocualitativo}{la misma muestra, dos preguntas}
Para determinar cuánto espacio ocupa una muestra interesa su volumen. Para
determinar cuánto principio activo contiene interesa su masa o su
concentración. La muestra es la misma, pero la pregunta cambia las magnitudes
que deben medirse y el modelo que resulta útil.
\end{ejemplocualitativo}

\begin{pregunta}
Una tableta se modela como una partícula para describir la posición de su centro
durante una caída. ¿Qué propiedad de la tableta se omite con ese modelo y en
qué situación dejaría de ser una omisión aceptable?
\end{pregunta}

\section{Magnitudes, unidades y dimensiones}

En mecánica, longitud, masa y tiempo funcionan como magnitudes fundamentales.
Sus símbolos dimensionales y unidades SI son:

\begin{center}
\begin{tabular}{lll}
\toprule
Magnitud & Dimensión & Unidad SI \\
\midrule
longitud & \(L\) & metro, \(\mathrm{m}\) \\
masa & \(M\) & kilogramo, \(\mathrm{kg}\) \\
tiempo & \(T\) & segundo, \(\mathrm{s}\) \\
\bottomrule
\end{tabular}
\end{center}

Una \textbf{dimensión} indica la naturaleza física de una cantidad. Una
\textbf{unidad} indica el patrón usado para medirla. Metro, centímetro y
kilómetro son unidades diferentes, pero las tres corresponden a la dimensión
longitud \(L\).

Las magnitudes derivadas se construyen combinando magnitudes fundamentales.
Si \(A\) representa área, \(V\) volumen, \(v\) rapidez y \(\rho\) densidad,
\[
  [A]=L^2,\qquad
  [V]=L^3,\qquad
  [v]=LT^{-1},\qquad
  [\rho]=ML^{-3}.
\]
Los corchetes se leen ``dimensión de''. Así, \([v]=LT^{-1}\) afirma que una
rapidez compara una longitud con un tiempo; no especifica si se mide en
\(\mathrm{m/s}\) o en \(\mathrm{km/h}\).

La densidad aparece cuando interesa comparar cuánta masa está contenida en un
volumen:
\[
  \rho\equiv\frac{m}{V}.
\]
Aquí \(m\) es la masa, \(V\) el volumen y \(\rho\) la densidad. En el SI,
\(\rho\) se expresa en \(\mathrm{kg/m^3}\). En laboratorio también es común
\(\mathrm{g/mL}\). Ambas unidades representan \(ML^{-3}\).

\begin{pregunta}
Dos muestras tienen densidades \(1{,}2\unidadSI{g/mL}\) y
\(1200\unidadSI{kg/m^3}\). ¿La segunda muestra es mil veces más densa o ambas
densidades describen el mismo valor físico?
\end{pregunta}

\section{Escalas, notación científica y prefijos}

La física compara tamaños muy distintos. Una molécula puede medirse en
nanómetros, una célula en micrómetros, una tableta en milímetros y una sala en
metros. La notación científica separa la cifra significativa de la potencia de
diez:
\[
  N=a\times10^n,
  \qquad 1\leq |a|<10.
\]

\begin{center}
\begin{tikzpicture}[x=1.25cm,y=1cm]
  \draw[-{Stealth[length=2mm]},thick] (0,0)--(8.4,0);
  \foreach \x/\e in {0/-9,2/-6,4/-3,6/0,8/3}{
    \draw (\x,0.12)--(\x,-0.12);
    \node[below,font=\small] at (\x,-0.12) {\(10^{\e}\unidadSI{m}\)};
  }
  \node[above,font=\small] at (0,0.15) {molécula};
  \node[above,font=\small] at (2,0.15) {célula};
  \node[above,font=\small] at (4,0.15) {tableta};
  \node[above,font=\small] at (6,0.15) {persona};
  \node[above,font=\small] at (8,0.15) {ciudad};
\end{tikzpicture}
\end{center}

Los prefijos sustituyen potencias de diez:
\[
\begin{array}{c@{\qquad}c@{\qquad}c}
\text{nano (n)}=10^{-9}
&
\text{micro (\(\mu\))}=10^{-6}
&
\text{mili (m)}=10^{-3}
\\[0.3em]
\text{centi (c)}=10^{-2}
&
\text{kilo (k)}=10^{3}.
\end{array}
\]
El prefijo forma parte de la unidad. Si la unidad se eleva a una potencia, el
factor del prefijo también:
\[
  1\unidadSI{cm^2}
  =(10^{-2}\unidadSI{m})^2
  =10^{-4}\unidadSI{m^2},
\]
\[
  1\unidadSI{cm^3}
  =(10^{-2}\unidadSI{m})^3
  =10^{-6}\unidadSI{m^3}.
\]
Esta última igualdad explica por qué
\(1\unidadSI{mL}=1\unidadSI{cm^3}=10^{-6}\unidadSI{m^3}\).

\section{Conversión de unidades}

Convertir una unidad no cambia la cantidad física; cambia su forma de
escritura. La igualdad
\[
  1\unidadSI{km}=10^3\unidadSI{m}
\]
permite construir dos factores iguales a uno:
\[
  \frac{10^3\unidadSI{m}}{1\unidadSI{km}}=1,
  \qquad
  \frac{1\unidadSI{km}}{10^3\unidadSI{m}}=1.
\]
Se escoge la orientación que cancela la unidad inicial.

Por ejemplo,
\[
72\frac{\mathrm{km}}{\mathrm{h}}
\left(\frac{10^3\unidadSI{m}}{1\unidadSI{km}}\right)
\left(\frac{1\unidadSI{h}}{3600\unidadSI{s}}\right)
=20\frac{\mathrm{m}}{\mathrm{s}}.
\]
El resultado expresa la misma rapidez: \(72\unidadSI{km/h}\) significa avanzar
\(20\unidadSI{m}\) cada segundo.

\Needspace{27\baselineskip}
\subsection{Ejemplo resuelto: masa de un excipiente líquido}

Se transfieren \(250\unidadSI{mL}\) de un excipiente líquido homogéneo cuya
densidad es
\[
  \rho=1{,}26\frac{\mathrm{g}}{\mathrm{mL}}.
\]
Se busca la masa transferida en kilogramos. El sistema es el líquido contenido
en la probeta. El modelo supone densidad uniforme, temperatura fija y ausencia
de pérdidas.

\begin{center}
\begin{tikzpicture}[scale=0.95]
  \draw[thick] (0,0)--(0.25,3.2)--(1.75,3.2)--(2,0)--cycle;
  \fill[AzulPregunta] (0.15,0.15)--(0.8,2.45)--(1.2,2.45)--(1.85,0.15)--cycle;
  \draw[AzulTitulo,thick] (0.76,2.45)--(1.24,2.45);
  \draw (2.15,2.45)--(2.45,2.45);
  \node[right,font=\small] at (2.45,2.45) {\(V=250\unidadSI{mL}\)};
  \node[font=\small] at (1,-0.35) {\(\rho=1{,}26\unidadSI{g/mL}\)};
\end{tikzpicture}
\end{center}

La densidad está definida por \(\rho=m/V\). Al despejar la masa,
\[
  m=\rho V.
\]
La sustitución conserva las unidades:
\begin{align*}
  m
  &=
  \left(1{,}26\frac{\mathrm{g}}{\mathrm{mL}}\right)
  (250\unidadSI{mL})
  =315\unidadSI{g},
  \\
  m
  &=
  315\unidadSI{g}
  \left(\frac{1\unidadSI{kg}}{10^3\unidadSI{g}}\right)
  =0{,}315\unidadSI{kg}.
\end{align*}
\resultado{\(m=0{,}315\unidadSI{kg}\).}

Las unidades de volumen se cancelan en la primera línea y los gramos en la
segunda. Además, como la densidad es algo mayor que
\(1\unidadSI{g/mL}\), la masa debe ser algo mayor que
\(250\unidadSI{g}\); \(315\unidadSI{g}\) es razonable.

La misma respuesta se obtiene en unidades SI:
\[
  \rho=1{,}26\times10^3\frac{\mathrm{kg}}{\mathrm{m^3}},
  \qquad
  V=2{,}50\times10^{-4}\unidadSI{m^3},
\]
\[
  m=\rho V
  =(1{,}26\times10^3)(2{,}50\times10^{-4})\unidadSI{kg}
  =0{,}315\unidadSI{kg}.
\]

\section{Análisis dimensional}

Una ecuación física solo puede igualar o sumar cantidades de la misma
dimensión. Esta exigencia se llama \textbf{homogeneidad dimensional}. En
\[
  x=x_0+vt,
\]
\(x\) y \(x_0\) son longitudes, mientras que
\[
  [vt]=(LT^{-1})T=L.
\]
Los tres términos tienen dimensión \(L\). En cambio, \(x=x_0+v+t\) no puede
representar una igualdad física: intenta sumar longitud, rapidez y tiempo.

La homogeneidad es una condición necesaria, pero no suficiente. Una expresión
dimensionalmente correcta todavía puede tener un signo, un factor numérico o
un supuesto equivocado.

\subsection{Forma posible de una distancia}

Supóngase que una distancia \(x\) depende únicamente de una aceleración
constante \(a\) y de un tiempo \(t\). Se propone
\[
  x=C\,a^n t^m,
\]
donde \(C\) no tiene dimensión. Como
\([x]=L\), \([a]=LT^{-2}\) y \([t]=T\),
\begin{align*}
  L
  &=(LT^{-2})^nT^m,
  \\
  L
  &=L^nT^{m-2n}.
\end{align*}
Las potencias de \(L\) y \(T\) deben coincidir en ambos lados:
\[
  n=1,\qquad m-2n=0,
\]
por lo que
\[
  n=1,\qquad m=2,\qquad x=C\,at^2.
\]
Las unidades confirman la forma:
\[
  \left(\frac{\mathrm{m}}{\mathrm{s^2}}\right)\mathrm{s^2}
  =\mathrm{m}.
\]
El análisis dimensional no determina \(C\). Tampoco demuestra que la posición
inicial y la velocidad inicial sean nulas. Cuando la cinemática incorpore esos
supuestos, aparecerá el caso particular
\(\Delta x=\tfrac12at^2\).

\begin{pregunta}
¿Cuál de las expresiones puede sumarse a una posición inicial \(x_0\)?
\[
  \text{A) }vt+at
  \qquad\qquad
  \text{B) }vt+at^2.
\]
La decisión debe justificarse comparando dimensiones término a término.
\end{pregunta}

\section{Estimaciones y orden de magnitud}

Una estimación busca la escala dominante de un resultado. Se declaran
supuestos simples, se redondean los datos y se operan potencias de diez. El
objetivo no es reemplazar una medición precisa, sino anticipar si una respuesta
debe estar en miles, millones o miles de millones.

\subsection{Ejemplo resuelto: respiraciones en un año}

Una frecuencia respiratoria media del orden de
\[
  f\sim10\frac{\text{respiraciones}}{\mathrm{min}}
\]
y la duración
\[
  1\ \text{año}
  \sim
  (60\ \text{min/h})(24\ \text{h/día})(365\ \text{días})
  \sim5\times10^5\ \text{min}
\]
producen
\[
  N
  \sim
  \left(10\frac{\text{resp}}{\text{min}}\right)
  (5\times10^5\text{ min})
  =5\times10^6\ \text{resp}.
\]
El orden de magnitud es \(10^7\) respiraciones por año. La frecuencia real
cambia con el sueño, el ejercicio y el estado de salud, pero la escala de
millones es robusta. Un resultado de \(10^3\) o \(10^{12}\) exigiría revisar
los datos, las unidades o las potencias de diez.

\begin{ejemplocualitativo}{una cifra que delata un error}
Para \(250\unidadSI{mL}\) de un líquido de densidad cercana a
\(1\unidadSI{g/mL}\), la masa debe estar cerca de cientos de gramos. Un
resultado de \(315\,000\unidadSI{kg}\) puede rechazarse antes de repetir la
aritmética: contradice la escala del sistema y convierte gramos a kilogramos en
la dirección equivocada.
\end{ejemplocualitativo}

\section{Síntesis}

\[
\begin{gathered}
\text{medición}
\longrightarrow
\text{número con unidad},
\\[0.4em]
\text{significado de la magnitud}
\longrightarrow
\text{dimensión},
\\[0.4em]
\text{ecuación}
\longrightarrow
\text{homogeneidad dimensional},
\\[0.4em]
\text{resultado}
\longrightarrow
\text{unidad, dimensión y escala razonable}.
\end{gathered}
\]

Las unidades no se agregan al final: participan en cada paso del cálculo. Las
dimensiones descartan relaciones imposibles y las estimaciones permiten
reconocer resultados absurdos. Para describir dónde está una partícula hará
falta, además, fijar un origen y una dirección positiva: ese será el paso desde
la medición hacia los sistemas de referencia y los vectores.

\end{document}
